\section{PHP}
    \par Considerando que o presente TCC tem como objetivo demonstrar a refatoração de uma API REST desenvolvida com o framework Laravel, o qual é baseado na linguagem de programação PHP, faz-se necessária uma fundamentação teórica sobre essa linguagem. Esta seção tem como propósito apresentar os principais conceitos do PHP antes de abordar o framework Laravel.

    \par A linguagem PHP, cujo nome é um acrônimo recursivo para PHP: Hypertext Preprocessor, teve um início singular, evoluindo de uma simples necessidade pessoal para uma das tecnologias mais difundidas na web. Sua história remonta a 1994, quando Rasmus Lerdorf desenvolveu uma série de scripts Interface Comum de Gateway (CGI, do inglês \textit{Common Gateway Interface}) para monitorar as visitas ao seu currículo online, um projeto que ele denominou ''Personal Home Page Tools'' \cite{livro:lockhart:2015:php}. 
    
    \par Esta primeira encarnação não era uma linguagem de script completa, mas sim um conjunto de ferramentas que ofereciam variáveis rudimentares e interpretação automática de formulários HTML. A grande virada ocorreu quando Andi Gutmans e Zeev Suraski reescreveram o núcleo da linguagem, transformando-a em uma linguagem de programação completa, com sintaxe mais consistente e suporte básico à orientação a objetos, culminando no lançamento do PHP 3 em 1998 \cite{livro:lockhart:2015:php}.
    
    \par A extensibilidade superior para diversos bancos de dados, protocolos e APIs foi um fator crucial que atraiu muitos desenvolvedores. Como resultado, ao final de 1998, o PHP 3 já estava instalado em cerca de 10 \% dos servidores web do mundo \cite{livro:lockhart:2015:php}.
    
    \par em relação à sua arquitetura, em sua essência, o PHP opera como uma linguagem de script interpretada e executada do lado do servidor (server-side). Seu modelo de funcionamento consiste em escrever o código, enviá-lo para um servidor web e executá-lo através de um interpretador a cada requisição. Tipicamente, é utilizado em conjunto com servidores como Apache ou Nginx para a geração de conteúdo dinâmico, embora também seja uma ferramenta poderosa para a construção de aplicações de linha de comando \cite{livro:lockhart:2015:php}.
    
    \par O processo de execução envolve a análise do código, sua compilação para um conjunto de instruções de máquina (Opcodes do Zend Engine) e, finalmente, a execução desse bytecode para gerar a resposta. Este ciclo, que ocorre a cada requisição, representa uma sobrecarga de processamento que soluções modernas, como caches de opcode, buscam mitigar para melhorar a performance da aplicação \cite{livro:lockhart:2015:php}.

    \par Tendo em vista a sua origem e sua arquitetura, em relação ao estado atual da linguagem PHP, considerando a data do presente trabalho, o PHP contemporâneo vive o que \cite{livro:lockhart:2015:php} descreve como uma "renascença", transformando-se em uma linguagem madura com funcionalidades avançadas como namespaces, traits e closures. Essa evolução não se restringiu à linguagem, mas impactou todo o seu ecossistema, que migrou de frameworks monolíticos para uma arquitetura baseada em componentes menores, especializados e interoperáveis \cite{livro:lockhart:2015:php}. 
    
    \par A introdução do gerenciador de dependências Composer mudou a forma como as aplicações são construídas, permitindo que desenvolvedores combinem componentes de diferentes fontes de maneira coesa. Essa interoperabilidade foi viabilizada pela criação de padrões comunitários propostos e mantidos pelo PHP Framework Interop Group (PHP-FIG) \cite{livro:lockhart:2015:php}.
    
    \par Consequentemente, as práticas de desenvolvimento também se modernizaram, substituindo métodos obsoletos como o envio de arquivos por FTP por fluxos de trabalho baseados em sistemas de controle de versão como o Git e o uso de ambientes de desenvolvimento virtualizados que espelham o servidor de produção \cite{livro:lockhart:2015:php}.

    \par Relacionado ao desenvolvimento e melhorias da linguagem PHP, o cenário atual é fortemente influenciado pela competição e pela busca por maior desempenho e padronização. Um marco importante foi a introdução do primeiro rascunho de uma especificação oficial da linguagem em 2014 \cite{livro:lockhart:2015:php}.